\documentclass[UTF8,a4paper,11pt]{ctexart}

\usepackage[margin=2.2cm]{geometry}
\usepackage{amsmath}
\usepackage{amssymb}
\usepackage{array}
\usepackage{booktabs}
\usepackage{longtable}
\usepackage{xcolor}
\usepackage{hyperref}

\hypersetup{
  colorlinks=true,
  linkcolor=blue!60!black,
  urlcolor=blue!60!black
}

\newcommand{\hex}[1]{\texttt{#1}}
\newcommand{\addr}[1]{\texttt{#1}}
\newcommand{\reg}[1]{\texttt{#1}}
\newcommand{\bit}[1]{\texttt{#1}}

\title{TX32F01 Cortex-M0 内存分配与外设寄存器地址图}
\author{Codex 生成，依据工程头文件与当前 Bootloader 布局整理}
\date{\today}

\begin{document}
\maketitle

\section{基本地址计算模型}

TX32F01 的外设采用 memory-mapped I/O。CPU 访问外设寄存器时，本质上是在访问某个固定地址。对任意外设 $P$，若其基地址为 $B_P$，寄存器偏移为 $O_R$，则寄存器绝对地址为
\[
  A(P,R) = B_P + O_R
\]

例如 UART 基地址为
\[
  B_{\mathrm{UART}} = \addr{0x49002000}
\]
UART 发送数据寄存器 \reg{TDR} 的偏移为
\[
  O_{\mathrm{TDR}} = \addr{0x10}
\]
因此：
\[
  A(\mathrm{UART},\mathrm{TDR})
  = \addr{0x49002000} + \addr{0x10}
  = \addr{0x49002010}
\]

对位域访问，若寄存器值为 $V$，字段起始 bit 为 $s$，字段宽度为 $w$，则字段 mask 为
\[
  M(s,w) = \left(2^w - 1\right) \ll s
\]
读取字段值：
\[
  F = \frac{V \mathbin{\&} M(s,w)}{2^s}
\]
写入字段值 $x$：
\[
  V' = \left(V \mathbin{\&} \sim M(s,w)\right)
       \mathbin{|}
       \left((x \ll s) \mathbin{\&} M(s,w)\right)
\]

例如 UART \reg{CER.UE} 位于 bit 0，宽度 1：
\[
  M_{\mathrm{UE}} = (2^1 - 1) \ll 0 = \addr{0x00000001}
\]
置位 UART 使能：
\[
  \reg{UART->CER} \gets \reg{UART->CER} \mathbin{|} \addr{0x00000001}
\]

\section{总体地址空间}

\begin{longtable}{>{\ttfamily}p{3.1cm}p{4.1cm}p{7cm}}
\toprule
地址范围或基址 & 区域 & 说明 \\
\midrule
\endhead
0x00000000 & Code alias / 启动映射区 & 复位后 Cortex-M0 从这里取初始 MSP 和 ResetHandler。实际内容映射到片内 Flash 起始区域。\\
0x01000000--0x01007FFF & Main Flash, 32 KB & 存放 Bootloader、APP、meta、boot flag。\\
0x20000000--0x20000FFF & SRAM, 4 KB & 存放软向量表、全局变量、栈、RTOS task stack、临时 buffer。\\
0x40001000 & SCU & 系统控制、时钟、外设使能、复位、低功耗。\\
0x40008000 & FLASH 控制器 & 控制 Flash 擦写、CRC、保护、Die ID。注意这不是代码存储地址。\\
0x49000000--0x4900B000 & 外设寄存器区 & TIM、SPI、GPIO、UART、I2C、ADC、IWDT、EXTI。\\
0xE000E000 起 & Cortex-M0 PPB & SysTick、NVIC、SCB 等 ARM 内核系统外设。\\
\bottomrule
\end{longtable}

\section{Flash 分区计算}

片内 Flash 总容量：
\[
  S_{\mathrm{flash}} = 32\,\mathrm{KB} = 32 \times 1024 = \addr{0x8000}
\]

Flash 起止地址：
\[
  F_{\mathrm{base}} = \addr{0x01000000}
\]
\[
  F_{\mathrm{end}}
  = F_{\mathrm{base}} + S_{\mathrm{flash}} - 1
  = \addr{0x01000000} + \addr{0x8000} - 1
  = \addr{0x01007FFF}
\]

当前工程采用如下分区：

\begin{longtable}{>{\ttfamily}p{3.3cm}>{\ttfamily}p{3.3cm}p{2.3cm}p{5.4cm}}
\toprule
起始地址 & 结束地址 & 大小 & 用途 \\
\midrule
\endhead
0x01000000 & 0x01001FFF & 8 KB & Bootloader：向量表、YMODEM、Flash 写入、CRC、跳转 APP。\\
0x01002000 & 0x010077FF & 22 KB & Application：APP/APP\_COOP/APP\_LOG/APP\_IRQ 等。\\
0x01007800 & 0x01007BFF & 1 KB & APP meta：magic、app\_size、crc16。\\
0x01007C00 & 0x01007FFF & 1 KB & Boot flag：如 \hex{0xA55AF00D} 表示强制留在 Bootloader。\\
\bottomrule
\end{longtable}

Bootloader 大小：
\[
  S_{\mathrm{BL}} = 8\,\mathrm{KB} = \addr{0x2000}
\]
\[
  F_{\mathrm{BL,end}}
  = \addr{0x01000000} + \addr{0x2000} - 1
  = \addr{0x01001FFF}
\]

APP 起始地址：
\[
  F_{\mathrm{APP,base}}
  = F_{\mathrm{base}} + S_{\mathrm{BL}}
  = \addr{0x01000000} + \addr{0x2000}
  = \addr{0x01002000}
\]

APP 最大容量：
\[
  S_{\mathrm{APP}}
  = \addr{0x01007800} - \addr{0x01002000}
  = \addr{0x5800}
  = 22\,\mathrm{KB}
\]

因此，OTA 接收 APP 镜像时必须满足：
\[
  0 < \mathrm{app\_size} \leq \addr{0x5800}
\]
写入地址范围必须满足：
\[
  \addr{0x01002000}
  \leq A_{\mathrm{write}}
  \leq
  \addr{0x010077FF}
\]

meta 区地址：
\[
  F_{\mathrm{meta,base}} = \addr{0x01007800}
\]
\[
  F_{\mathrm{meta,end}} = \addr{0x01007BFF}
\]

boot flag 区地址：
\[
  F_{\mathrm{flag,base}} = \addr{0x01007C00}
\]
\[
  F_{\mathrm{flag,end}} = \addr{0x01007FFF}
\]

\section{SRAM 分配计算}

SRAM 总容量：
\[
  S_{\mathrm{sram}} = 4\,\mathrm{KB} = \addr{0x1000}
\]

SRAM 起止地址：
\[
  R_{\mathrm{base}} = \addr{0x20000000}
\]
\[
  R_{\mathrm{end}}
  = R_{\mathrm{base}} + S_{\mathrm{sram}} - 1
  = \addr{0x20000000} + \addr{0x1000} - 1
  = \addr{0x20000FFF}
\]

因为 Cortex-M0 没有 VTOR，当前 Bootloader 设计在 SRAM 头部保留软向量表。软向量表大小为 96 B：
\[
  S_{\mathrm{softvec}} = 96 = \addr{0x60}
\]

软向量表范围：
\[
  R_{\mathrm{softvec,base}} = \addr{0x20000000}
\]
\[
  R_{\mathrm{softvec,end}}
  = \addr{0x20000000} + \addr{0x60} - 1
  = \addr{0x2000005F}
\]

APP 可用 RAM 起始地址必须避开软向量表：
\[
  R_{\mathrm{APP,base}}
  = R_{\mathrm{softvec,end}} + 1
  = \addr{0x20000060}
\]

APP 最大可用 SRAM：
\[
  S_{\mathrm{APP,RAM}}
  = \addr{0x20001000} - \addr{0x20000060}
  = \addr{0xFA0}
  = 4000\,\mathrm{B}
\]

实际 APP 还必须在这 4000 B 中同时容纳：
\[
  S_{\mathrm{RW}} + S_{\mathrm{ZI}} + S_{\mathrm{heap}}
  + S_{\mathrm{MSP}} + \sum_i S_{\mathrm{PSP},i}
  \leq 4000\,\mathrm{B}
\]

裸机或协作式调度器通常只有一个主栈：
\[
  S_{\mathrm{RW}} + S_{\mathrm{ZI}} + S_{\mathrm{buffer}} + S_{\mathrm{MSP}}
  \leq 4000\,\mathrm{B}
\]

FreeRTOS 需要额外考虑每个任务的独立栈：
\[
  S_{\mathrm{RTOS}}
  =
  S_{\mathrm{kernel}}
  + S_{\mathrm{idle\_stack}}
  + \sum_{i=1}^{N} S_{\mathrm{task},i}
  + S_{\mathrm{queues}}
  + S_{\mathrm{timers}}
\]

\section{外设基地址与地址公式}

\begin{longtable}{p{2.3cm}>{\ttfamily}p{3.2cm}p{7.7cm}}
\toprule
模块 & 基地址 & 地址计算 \\
\midrule
\endhead
SCU & 0x40001000 & $A = \addr{0x40001000} + O$ \\
FLASH 控制器 & 0x40008000 & $A = \addr{0x40008000} + O$ \\
TIM0 & 0x49000000 & $A = \addr{0x49000000} + O$ \\
TIM1 & 0x49000400 & $A = \addr{0x49000400} + O$ \\
TIM2 & 0x49000800 & $A = \addr{0x49000800} + O$ \\
SPI & 0x49000C00 & $A = \addr{0x49000C00} + O$ \\
GPIO0 & 0x49001000 & $A = \addr{0x49001000} + O$ \\
GPIO1 & 0x49001400 & $A = \addr{0x49001400} + O$ \\
GPIO2 & 0x49001800 & $A = \addr{0x49001800} + O$ \\
GPIO3 & 0x49001C00 & $A = \addr{0x49001C00} + O$ \\
UART & 0x49002000 & $A = \addr{0x49002000} + O$ \\
I2C & 0x49005000 & $A = \addr{0x49005000} + O$ \\
ADC & 0x49007000 & $A = \addr{0x49007000} + O$ \\
IWDT & 0x49009000 & $A = \addr{0x49009000} + O$ \\
EXTI & 0x4900B000 & $A = \addr{0x4900B000} + O$ \\
\bottomrule
\end{longtable}

Timer 和 GPIO 都是重复实例。Timer 的第 $n$ 个实例可以写为：
\[
  B_{\mathrm{TIM}n} = \addr{0x49000000} + n \times \addr{0x400},
  \quad n \in \{0,1,2\}
\]

GPIO 的第 $n$ 个实例可以写为：
\[
  B_{\mathrm{GPIO}n} = \addr{0x49001000} + n \times \addr{0x400},
  \quad n \in \{0,1,2,3\}
\]

例如 GPIO3 的输出数据寄存器 \reg{ODR} 偏移为 \addr{0x28}：
\[
  A(\mathrm{GPIO3},\mathrm{ODR})
  =
  \left(\addr{0x49001000} + 3 \times \addr{0x400}\right)
  + \addr{0x28}
  =
  \addr{0x49001C28}
\]

\section{SCU 寄存器}

SCU 基地址：
\[
  B_{\mathrm{SCU}} = \addr{0x40001000}
\]

\begin{longtable}{>{\ttfamily}p{2cm}>{\ttfamily}p{3cm}p{2.6cm}p{6cm}}
\toprule
偏移 & 绝对地址 & 寄存器 & 说明 \\
\midrule
\endhead
0x00 & 0x40001000 & PKR & 访问保护 key。\\
0x04 & 0x40001004 & CCR & HCLK 分频、MCO 输出。\\
0x08 & 0x40001008 & PENR & 外设时钟使能。\\
0x0C & 0x4000100C & PRSTR & 外设复位。\\
0x10 & 0x40001010 & RSR & 复位原因。\\
0x14 & 0x40001014 & PMCR & BOR/PVD 电源监控。\\
0x18 & 0x40001018 & LPCR & 低功耗配置。\\
0x1C & 0x4000101C & DBGCR & debug sleep/stop、Timer/IWDT 调试暂停。\\
0x20 & 0x40001020 & STR & SysTick 校准值。\\
0x24 & 0x40001024 & BSR & Boot 映射选择。\\
\bottomrule
\end{longtable}

外设时钟使能的一般形式：
\[
  \reg{SCU->PENR} \gets \reg{SCU->PENR} \mathbin{|} (1 \ll k)
\]
其中 $k$ 为对应外设的 enable bit，例如 UART 的时钟使能位为 bit 3：
\[
  \reg{SCU->PENR} \gets \reg{SCU->PENR} \mathbin{|} (1 \ll 3)
\]

\section{FLASH 控制器寄存器}

FLASH 控制器基地址：
\[
  B_{\mathrm{FLASHCTL}} = \addr{0x40008000}
\]

\begin{longtable}{>{\ttfamily}p{2cm}>{\ttfamily}p{3cm}p{2.6cm}p{6cm}}
\toprule
偏移 & 绝对地址 & 寄存器 & 说明 \\
\midrule
\endhead
0x00 & 0x40008000 & ASR & Flash 访问状态。\\
0x04 & 0x40008004 & ACR & 访问控制、wait cycle。\\
0x08 & 0x40008008 & ACHKR & 擦写校验使能。\\
0x0C & 0x4000800C & PEKEYR & Program/erase key。\\
0x20 & 0x40008020 & IFR & Flash 状态/中断标志。\\
0x24 & 0x40008024 & IER & Flash 中断使能。\\
0x30 & 0x40008030 & TKEYR & timing 配置 key。\\
0x34 & 0x40008034 & TNVS & Flash timing。\\
0x38 & 0x40008038 & TPROG & Program timing。\\
0x3C & 0x4000803C & TADS & Address setup timing。\\
0x40 & 0x40008040 & TADH & Address hold timing。\\
0x44 & 0x40008044 & TPGS & Program setup timing。\\
0x48 & 0x40008048 & TPGH & Program hold timing。\\
0x4C & 0x4000804C & TPRCV & Program recovery timing。\\
0x50 & 0x40008050 & TSRCV & Sector erase recovery timing。\\
0x54 & 0x40008054 & TCRCV & Chip erase recovery timing。\\
0x58 & 0x40008058 & TSERASE & Sector erase timing。\\
0x5C & 0x4000805C & TCERASE & Chip erase timing。\\
0x60 & 0x40008060 & TRW & Read/write timing。\\
0xE0 & 0x400080E0 & DIEID0 & Die ID 0。\\
0xE4 & 0x400080E4 & DIEID1 & Die ID 1。\\
0xE8 & 0x400080E8 & DIEID2 & Die ID 2。\\
0xEC & 0x400080EC & DIEID3 & Die ID 3。\\
0xF0 & 0x400080F0 & CRCCR & CRC 控制。\\
0xF4 & 0x400080F4 & CRCARL & CRC 地址低寄存器。\\
0xF8 & 0x400080F8 & CRCARH & CRC 地址高寄存器。\\
0xFC & 0x400080FC & CRCVR & CRC 结果。\\
0x100 & 0x40008100 & PKEYR & 控制寄存器访问保护 key。\\
0x104 & 0x40008104 & PRGKEYR & Main Flash program enable key。\\
0x10C & 0x4000810C & NVR\_WR\_EN & NVR 写使能。\\
\bottomrule
\end{longtable}

Flash CRC 地址范围要特别注意。若硬件 API 采用 inclusive end，则对长度为 $N$ 字节的 APP：
\[
  A_{\mathrm{crc,start}} = F_{\mathrm{APP,base}}
\]
\[
  A_{\mathrm{crc,end}} = F_{\mathrm{APP,base}} + N - 1
\]
而不是 $F_{\mathrm{APP,base}} + N$。

\section{Timer 寄存器}

Timer 基地址：
\[
  B_{\mathrm{TIM}n} = \addr{0x49000000} + n \times \addr{0x400}
\]

\begin{longtable}{>{\ttfamily}p{2cm}p{3cm}p{8cm}}
\toprule
偏移 & 寄存器 & 说明 \\
\midrule
\endhead
0x00 & CR & 计数器使能、运行模式、预装载。\\
0x04 & CCCR & Capture/Compare 控制。\\
0x08 & BCR & Break 控制。\\
0x0C & TCR & Trigger 控制。\\
0x10 & IER & 中断使能。\\
0x14 & IFR & 中断标志。\\
0x18 & CNT & 当前计数值。\\
0x1C & DIV & 分频。\\
0x20 & ARR & 自动重装载值。\\
0x24 & CCR & Compare/Capture 值。\\
0x28 & DTG & 死区时间。\\
\bottomrule
\end{longtable}

PWM 频率公式可以写为：
\[
  f_{\mathrm{PWM}}
  =
  \frac{f_{\mathrm{PCLK}}}{(DIV+1)(ARR+1)}
\]

若系统时钟为 $24\,\mathrm{MHz}$，分频寄存器选择 $DIV=23$，自动重装载 $ARR=999$：
\[
  f_{\mathrm{PWM}}
  =
  \frac{24\,000\,000}{(23+1)(999+1)}
  =
  1000\,\mathrm{Hz}
\]

PWM 占空比近似为：
\[
  D \approx \frac{CCR}{ARR+1}
\]

\section{GPIO 寄存器}

GPIO 基地址：
\[
  B_{\mathrm{GPIO}n} = \addr{0x49001000} + n \times \addr{0x400}
\]

\begin{longtable}{>{\ttfamily}p{2cm}p{3cm}p{8cm}}
\toprule
偏移 & 寄存器 & 说明 \\
\midrule
\endhead
0x00 & MDR & GPIO 模式。\\
0x04 & PUR & 上拉。\\
0x08 & PDR & 下拉。\\
0x0C & DSR & 驱动强度。\\
0x10 & SRR & Slew rate。\\
0x14 & OTR & 推挽/开漏。\\
0x18 & IER & 输入使能。\\
0x20 & AFR & 复用功能选择。\\
0x28 & ODR & 输出数据。\\
0x2C & BSR & 置位输出。\\
0x30 & BRR & 清零输出。\\
0x34 & BFR & 翻转输出。\\
0x38 & IDR & 输入数据。\\
\bottomrule
\end{longtable}

对 GPIO 第 $p$ 个 pin，置位输出通常使用：
\[
  \reg{GPIOx->BSR} \gets 1 \ll p
\]
清零输出：
\[
  \reg{GPIOx->BRR} \gets 1 \ll p
\]
翻转输出：
\[
  \reg{GPIOx->BFR} \gets 1 \ll p
\]

相对于读改写 \reg{ODR}，\reg{BSR}/\reg{BRR}/\reg{BFR} 更适合中断和并发场景。

\section{UART 寄存器}

UART 基地址：
\[
  B_{\mathrm{UART}} = \addr{0x49002000}
\]

\begin{longtable}{>{\ttfamily}p{2cm}>{\ttfamily}p{3cm}p{2.5cm}p{6cm}}
\toprule
偏移 & 绝对地址 & 寄存器 & 说明 \\
\midrule
\endhead
0x00 & 0x49002000 & CR & 校验、字长、停止位、极性。\\
0x04 & 0x49002004 & BRR & 波特率分频。\\
0x08 & 0x49002008 & IER & 中断使能。\\
0x0C & 0x4900200C & ISR & 状态/中断标志。\\
0x10 & 0x49002010 & TDR & 发送数据。\\
0x14 & 0x49002014 & RDR & 接收数据。\\
0x20 & 0x49002020 & CER & UART/TX/RX 使能。\\
0x100 & 0x49002100 & SADJ & stop bit 宽度微调。\\
\bottomrule
\end{longtable}

UART 发送寄存器地址计算：
\[
  A_{\mathrm{UART,TDR}}
  =
  \addr{0x49002000} + \addr{0x10}
  =
  \addr{0x49002010}
\]

UART 接收寄存器地址计算：
\[
  A_{\mathrm{UART,RDR}}
  =
  \addr{0x49002000} + \addr{0x14}
  =
  \addr{0x49002014}
\]

当前工程常用 UART 配置为 115200 8N1。若发送一个字节包含 1 start、8 data、1 stop，则一字节线路时间为：
\[
  t_{\mathrm{byte}}
  =
  \frac{10}{115200}
  \approx 86.8\,\mu s
\]

这说明 polling UART 单字节发送可以接受，但长字符串打印会明显占用 CPU：
\[
  t_{\mathrm{print}} \approx N_{\mathrm{bytes}} \times 86.8\,\mu s
\]

\section{SPI、I2C、ADC、IWDT、EXTI 寄存器摘要}

\subsection{SPI}

\[
  B_{\mathrm{SPI}} = \addr{0x49000C00}
\]

\begin{center}
\begin{tabular}{>{\ttfamily}p{2cm}p{3cm}p{8cm}}
\toprule
偏移 & 寄存器 & 说明 \\
\midrule
0x00 & CR & SPI 使能、主从、CPOL/CPHA、NSS、FIFO 清除。\\
0x04 & SR & 传输状态、FIFO 状态、错误标志。\\
0x08 & IER & SPI 中断使能。\\
0x0C & DIVR & SPI 分频。\\
0x10 & DATAR & SPI 数据。\\
0x14 & TCNTR & TX FIFO 计数。\\
0x18 & RCNTR & RX FIFO 计数。\\
\bottomrule
\end{tabular}
\end{center}

\subsection{I2C}

\[
  B_{\mathrm{I2C}} = \addr{0x49005000}
\]

\begin{center}
\begin{tabular}{>{\ttfamily}p{2cm}p{3cm}p{8cm}}
\toprule
偏移 & 寄存器 & 说明 \\
\midrule
0x00 & DR & 收发数据。\\
0x04 & AR & 从机地址、地址 mask、General Call。\\
0x08 & CR & START、STOP、ACK、使能、控制码。\\
0x0C & SR & I2C 状态码。\\
\bottomrule
\end{tabular}
\end{center}

\subsection{ADC}

\[
  B_{\mathrm{ADC}} = \addr{0x49007000}
\]

\begin{longtable}{>{\ttfamily}p{2cm}>{\ttfamily}p{3cm}p{2.5cm}p{6cm}}
\toprule
偏移 & 绝对地址 & 寄存器 & 说明 \\
\midrule
\endhead
0x00 & 0x49007000 & IFR & ADC 中断标志。\\
0x04 & 0x49007004 & IER & ADC 中断使能。\\
0x08 & 0x49007008 & CR & ADC 启动/使能。\\
0x0C & 0x4900700C & CFGR & clock、trigger、VREF。\\
0x10 & 0x49007010 & SQR1 & 序列控制 1。\\
0x14 & 0x49007014 & SQR2 & 序列控制 2。\\
0x20 & 0x49007020 & DR1 & 序列 1 结果。\\
0x24 & 0x49007024 & DR2 & 序列 2 结果。\\
0x28 & 0x49007028 & DR3 & 序列 3 结果。\\
0x2C & 0x4900702C & DR4 & 序列 4 结果。\\
0x30 & 0x49007030 & DR5 & 序列 5 结果。\\
0x34 & 0x49007034 & DR6 & 序列 6 结果。\\
0x38 & 0x49007038 & DR7 & 序列 7 结果。\\
0x3C & 0x4900703C & DR8 & 序列 8 结果。\\
0x40 & 0x49007040 & OFR & offset。\\
0x44 & 0x49007044 & VREFRL & 2.048 V 内部参考实测值。\\
0x48 & 0x49007048 & VREFRH & 高参考实测值。\\
\bottomrule
\end{longtable}

若 ADC 为 $N$ bit，原始码值为 $C$，参考电压为 $V_{\mathrm{ref}}$，则输入电压近似为：
\[
  V_{\mathrm{in}} \approx \frac{C}{2^N - 1} V_{\mathrm{ref}}
\]

若用内部参考反推 VDD，内部参考实际电压为 $V_{\mathrm{intref}}$，ADC 测得码值为 $C_{\mathrm{intref}}$：
\[
  V_{\mathrm{DD}}
  \approx
  \frac{(2^N - 1)V_{\mathrm{intref}}}{C_{\mathrm{intref}}}
\]

\subsection{IWDT}

\[
  B_{\mathrm{IWDT}} = \addr{0x49009000}
\]

\begin{center}
\begin{tabular}{>{\ttfamily}p{2cm}p{3cm}p{8cm}}
\toprule
偏移 & 寄存器 & 说明 \\
\midrule
0x00 & KR & 看门狗 key。\\
0x04 & PR & 预分频。\\
0x08 & RLR & 重装载值。\\
0x0C & SR & 状态。\\
0x10 & IER & 中断使能。\\
0x14 & IFR & 中断标志。\\
\bottomrule
\end{tabular}
\end{center}

\subsection{EXTI}

\[
  B_{\mathrm{EXTI}} = \addr{0x4900B000}
\]

\begin{center}
\begin{tabular}{>{\ttfamily}p{2cm}p{3cm}p{8cm}}
\toprule
偏移 & 寄存器 & 说明 \\
\midrule
0x00 & IMR & 中断 mask。\\
0x04 & EMR & 事件 mask。\\
0x08 & RTSR & 上升沿触发选择。\\
0x0C & FTSR & 下降沿触发选择。\\
0x10 & SWIER & 软件触发。\\
0x14 & PR & pending 标志。\\
0x18 & CFGR & EXTI line 到 GPIO port 映射。\\
\bottomrule
\end{tabular}
\end{center}

\section{Cortex-M0 系统外设区}

Cortex-M0 标准 PPB 区域常用地址如下：

\begin{longtable}{p{3cm}>{\ttfamily}p{3.6cm}p{7cm}}
\toprule
模块 & 地址 & 说明 \\
\midrule
\endhead
SysTick & 0xE000E010 & \reg{CTRL}、\reg{LOAD}、\reg{VAL}、\reg{CALIB}。\\
NVIC & 0xE000E100 & IRQ enable、disable、pending、priority。\\
SCB & 0xE000ED00 & CPUID、ICSR、AIRCR、SCR、CCR、SHPR、SHCSR。\\
\bottomrule
\end{longtable}

当前工程特别要注意：
\[
  \text{Cortex-M0 has no VTOR}
\]
因此不能使用 \reg{SCB->VTOR} 改变向量表基址。APP 中断必须通过 Bootloader 的硬件向量表和 SRAM 软向量表转发。

\section{使用建议}

\begin{enumerate}
  \item 写外设驱动前，先用 \reg{SCU->PENR} 打开对应外设时钟。
  \item 修改外设配置前，必要时用 \reg{SCU->PRSTR} 复位该外设。
  \item APP 的 RAM 起始地址必须从 \addr{0x20000060} 开始，避开软向量表。
  \item Flash 写入必须限制在 APP 区：\addr{0x01002000}--\addr{0x010077FF}。
  \item Flash 控制器地址 \addr{0x40008000} 是寄存器区，不是代码 Flash 存储区。
  \item GPIO 置位/清零/翻转优先使用 \reg{BSR}/\reg{BRR}/\reg{BFR}，避免读改写竞争。
  \item 4 KB SRAM 极小，应定期检查 map 文件，特别是 task stack、printf buffer 和日志 page buffer。
\end{enumerate}

\end{document}
